\documentclass[10pt]{article}

% Two builds from one source, as TMLR requires: the review build (default) is
% anonymized and headed "Under review as submission to TMLR"; the preprint build
% (\def\TMLRSUBMISSION{1}) prints the author block for arXiv.  Only the style
% option changes, so the text is identical between the two.
\ifdefined\TMLRSUBMISSION
  \usepackage[preprint]{tmlr}
\else
  \usepackage{tmlr}
\fi
\usepackage[utf8]{inputenc}
\usepackage[T1]{fontenc}
\usepackage{hyperref}
\usepackage{url}
\usepackage{booktabs}
\usepackage{amsfonts}
\usepackage{amsmath}
\usepackage{microtype}
\usepackage{graphicx}
\usepackage{xcolor}
\usepackage{underscore}

\title{Sharing a Route Is Not Free: Measuring the Cost of Cross-Layer Top-$k$ Reuse in Real Weights}

\author{%
  Anonymous Author(s) \\
  Affiliation \\
  Address \\
  \texttt{email}
}

\def\month{MM}
\def\year{YYYY}
\def\openreview{\url{https://openreview.net/forum?id=XXXX}}

\begin{document}

\maketitle

\begin{abstract}
Families of sparse attention let one layer reuse the token selection another
layer computed, resting on adjacent layers choosing nearly the same tokens -- a claim about \emph{selection}, not \emph{quality}.
IndexCache \citep{indexcache} measures quality at benchmark granularity on
trained sparse models; on a checkpoint trained with sharing (an FSS route
pattern) the group's own set costs twelve times less than an arbitrary
substitute of the same size.
We isolate the selection-set variable on dense pretrained weights with a
three-readout design separating token-level from decision-level cost, patching
the attention pass so each layer keeps its own queries, keys and values and only
the attended set is borrowed; dense, per-layer top-$k$ and group-leader reuse
share weights, tokens and batch. The harness refuses to report unless per-layer
selection at $k$ equal to the sequence length matches the dense pass; that
control caught a real defect. On
twenty batches of 512 positions the \emph{mean} cost per borrowed pair is 11.9
times higher when the source is layer 0 or 1 than 2 or later. Shifting groups
by one layer cuts cost 2--3.4 times at three of four widths. Three similarity-shaped predictors fail to rank the cost: overlap $r=-0.357$ on
23 Qwen2.5-0.5B pairs, CKA $r=-0.341$ and contribution $r=-0.071$ on a GPT-2
probe (\S\ref{sec:whatnot}). Against an
arbitrary $k$-set of the same size the reuse cost is $0.034$ to $0.126$ of that
null: most of it is the price of a set a layer did not choose. At the decision level on ARC-Easy, 0.5B loses at most
25.5 points across 25 cells, five of them at $p<0.05$; at $k$ equal to 67\%
of the decision-level context (256 of 384 positions) the largest difference is
1 point. Reuse destroys more answers than it repairs (10.1 versus 7.9 per hundred
at-risk on ARC-Easy, 12.4 versus 6.9 on ARC-Challenge), an asymmetry of 1.3--1.8
times rather than 5:1, and exemplar count alone can flip its sign.
\end{abstract}

\section{Introduction}

Sparse attention makes long context affordable by attending to a subset of
positions. A recent line of work notices that the subsets chosen by
neighbouring layers are similar enough to share, and uses that to avoid
recomputing an index per layer. Multi-layer Indexer reports that adjacent layers
select almost the same tokens and that the sharing needs no retraining
\citep{mli}. CLSA shares both the key-value blocks and the routing across layers
\citep{clsa}. IndexCache reuses indices across layers as an engineering
optimisation \citep{indexcache}.

The empirical premise in these three papers is a measurement of \emph{overlap}: the
fraction of positions that two layers both select. Our own reproduction
confirms that premise is real. Adjacent layers overlap far above chance,
and the enrichment grows with sequence length (1.19 times chance at 512
positions with $k=128$, 1.96 times at 2048 positions).

What the premise does not establish is that the borrowed selection still
attends to the right things. Two layers can select the same positions for
different reasons: layer $\ell{+}1$ may select position $j$ because its own
query needs it, while layer $\ell$ selected it for a reason that does not
transfer. If so, overlap is high and reuse is still harmful. IndexCache
\citep{indexcache} measures exactly this quality loss, but at benchmark
granularity on trained sparse models and by retention ratio; Kascade
\citep{kascade} uses a related mass-overlap ratio as a systems heuristic. We
isolate the same question at a different grain: on dense pretrained weights we
hold queries, keys, values and batch fixed and change only the set of attended
positions, then read the cost at both token and decision level. A systematic
study of trained \emph{key-value} sharing \citep{kvshare} answers a different
question because it retrains; the throughput-oriented papers
\citep{mli,clsa} report latency, not this quality term.

This paper measures the step. We take a pretrained model, patch every attention
layer so that the set of attended positions can be supplied from outside, and
compare three readings that differ in nothing else:

\begin{itemize}
\item \textbf{dense}: every position attended, the model as shipped.
\item \textbf{own}: each layer attends to its own top-$k$ positions.
\item \textbf{reuse}: each layer attends to the top-$k$ positions of the
      first layer of its group.
\end{itemize}

All three use each layer's own queries, keys and values; only the \emph{set}
differs. We report $\text{reuse} - \text{own}$, which removes the cost of
sparsification itself and isolates the cost of sharing.

\paragraph{What we find.}
(i) Reuse costs quality at the token level, and the cost is governed by $k$ and
by the identity of the group leader, and by inter-layer distance much less.
(ii) The measurement that motivates reuse (index overlap) is only a weak
predictor of cost, and part of what it appears to explain is confounded with
leader identity; a mass-based diagnostic improves on it in a head-to-head test
but is not an identified mechanism.
(iii) At the decision level the cost is much smaller than the token-level
number suggests, and in most configurations not statistically detectable at
our sample sizes.
(iv) The direction of the residual damage is observed consistently across the
significant cells but is modest, and our first estimate of it was wrong for an
algebraic reason that we now state as a gate.

\paragraph{What this paper is not.}
It is not a claim that reuse is useless. At $k$ equal to 75\% of the sequence
and groups of up to four layers, the cost is under 0.0122 nats of token
NLL, and at the decision level we find configurations where reuse is slightly
\emph{better} than per-layer selection. What it is not is a measurement of a
model trained to \emph{share}: outside \S\ref{sec:leader} every number borrows
routing post hoc from weights that received no training signal for reuse. The
checkpoints we run include one that learned to select per layer and one trained
with sharing, and what the second shows is narrow: substituting a random set for
the group's own set costs twelve times more, so training does not make a layer
indifferent to which set it receives. A system trained with sharing could still
move its representations to make sharing cheaper in ways a fixed-weight
substitution does not see.

\section{Instrumentation}

\subsection{Three readings, one model}

We patch each attention module's forward pass. The patch computes attention
scores exactly as the original does, then masks all but a supplied set of
positions per query, applies the softmax over the masked row, and multiplies
by the layer's own values. The dense reading supplies the full triangular mask
and must reproduce the unpatched model bit for bit; the other two supply
per-layer or inherited selections.

Three implementation details cost us real time and are worth stating because
each produced plausible-looking wrong numbers before it was fixed.

\begin{enumerate}
\item \textbf{Axis.} Scores have shape (batch, head, query, \emph{key}). Selections
      therefore live on the last axis. A gather on the query axis returns a
      tensor of the right shape with the wrong contents and does not raise.
\item \textbf{Masking, not gathering.} Extracting the kept scores and scattering
      them back changes the order of the softmax reduction. In bfloat16,
      addition is not associative, and the result is no longer bit-identical
      to the dense pass even when $k$ equals the sequence length. We therefore
      mask: keep a boolean of shape (batch, head, query, key), set the kept
      entries, and let the softmax reduce over the full row.
\item \textbf{One contiguous range.} A triangular mask must not be added to a mask
      built from \texttt{-inf}: the sum of two infinities is not a number. We
      use the dtype's most negative finite value.
\end{enumerate}

\subsection{The control that must pass}

Before any reading is reported, the instrument runs the per-layer reading with
$k$ equal to the sequence length and asserts that its mean token NLL equals the
dense value to within 0.000001. This is not a formality: the wrong-axis
implementation above satisfied every shape check and produced a loss far outside
that tolerance, which the control rejected. The rejections are kept as artifacts
rather than narrated: \texttt{paper/evidence/run-logs/main.log} is a failed
control at the tolerance of the loss, and
\texttt{real/data/control\_failure\_7b.txt} is a failed one
at the boundary of the argument. The control value is recorded in every
artifact, and a failed control invalidates the cell rather than the run.

Threshold on the control is deliberately tight, and this is a design choice
rather than an accident: a loose tolerance would have let the bfloat16
reduction-order error pass as noise, and that error is one of the mechanisms
by which the measurement could silently become a measurement of the harness.

\subsection{Models, data, benchmark}

We use Qwen2.5-0.5B (24 layers, hidden 896, 14 heads, 2 key-value heads),
Qwen2.5-1.5B and Qwen2.5-7B (28 layers each) in bfloat16, on wikitext-2-raw-v1
test for the token-level measurements and on ARC-Easy and ARC-Challenge for the
decision-level ones. All runs execute on a single AMD Radeon instance with ROCm.
Token-level grids are measured at twenty batches and single-pair sweeps at
twenty batches as well; the four-batch run reported in an earlier version of
this work is kept as a sample-size control, not as a result.

We report two quantities throughout. Token-level: mean per-token
negative log-likelihood, and its difference from the per-layer reading. Decision-level:
exact accuracy over four-way multiple choice, with each option scored by the sum
of its tokens' log-probabilities, without decoding, because a decoding
policy would confound the attribution.

\section{Token-level cost}

\subsection{The grid}

Table~\ref{tab:grid} gives $\text{reuse} - \text{own}$ for 0.5B at 512 positions.
Sparsification alone is nearly free in this range and sometimes slightly
helpful ($k=64$ to $k=256$ \emph{lower} NLL, by at most 0.0016), so every cost
below belongs to sharing rather than to sparsity. At $k=32$ sparsification
itself costs 0.0107, which is why that row is read as a sparsity effect on
top of a sharing cost.

\begin{table}[h]
\centering
\caption{$\text{reuse}-\text{own}$ in nats per token, Qwen2.5-0.5B, 512
positions, twenty batches, dense 2.82874 (perplexity 16.92). Control exact
($0$) in every cell. Source: \texttt{real/data/main\_20b\_05.json}.}
\label{tab:grid}
\begin{tabular}{rrrrrr}
\toprule
$k$ & $g{=}2$ & $g{=}3$ & $g{=}4$ & $g{=}6$ & $g{=}12$ \\
\midrule
32  & +0.8050 & +0.4854 & +0.9615 & +1.4553 & +3.2594 \\
64  & +0.4447 & +0.1884 & +0.5258 & +0.8356 & +2.4429 \\
128 & +0.1726 & +0.0602 & +0.2302 & +0.3581 & +1.6102 \\
192 & +0.0678 & +0.0292 & +0.1169 & +0.1731 & +0.9188 \\
256 & +0.0242 & +0.0114 & +0.0560 & +0.0736 & +0.4493 \\
384 & +0.0016 & +0.0004 & +0.0122 & +0.0154 & +0.0639 \\
\bottomrule
\end{tabular}
\end{table}

Three regularities. First, $k$ is the main switch, and the useful range is
narrow: at $k=384$ (75\% of the sequence) groups up to four layers cost at
most $+0.0122$ nats, taking perplexity from 16.92 to 17.13; at $k=32$ the
same grouping costs between $+0.8050$ and $+3.2594$, which is perplexity 37.9
to 440.6.

Second, cost is not monotone in group width, and the violation is stable
across samples: $g{=}3$ is \emph{cheaper} than $g{=}2$ at every $k$
($+0.0016$ against $+0.0004$ at $k=384$, $+0.1726$ against $+0.0602$ at
$k=128$), while $g{=}4$ is dearer than both. The reason is not that wide
groups help: a three-layer grouping does not ask layer 3 to borrow, because
layer 3 leads its own group, and layer 3 is the most expensive
\emph{follower} we measured. In an even grouping layer 3 follows layer 2 and
pays $+0.0146$; \S\ref{sec:leader} separates leader from follower.

Third, larger models share more cheaply at the same $k$: $k{=}128$, $g{=}2$
costs $+0.1726$ nats on 0.5B, $+0.0396$ on 1.5B and $+0.0243$ on 7B, factors
of 0.23 and 0.14. This is the one regularity that held identically in the
four-batch sample.

\subsection{What the reuse cost is measured against}

Every cell above compares a follower using the \emph{leader's} set against the
same layer using its \emph{own} set, so the two states differ twice over: the
set belongs to someone else, and it is not the set this layer would have
chosen. Those two differences can be separated, because the leader's set is
something we can replace with an arbitrary set of the same size. We do that at
each $(k,g)$ cell: the leaders keep the selection they would have made, every
borrowing layer receives a random $k$-set drawn only from positions it can see, and the null
perturbs exactly the layers the reuse cell perturbs. On 0.5B at $k{=}256$ the
null costs $+0.31493$, $+0.33666$ and $+0.58492$ nats for $g{=}2,3,6$, and on
7B $+0.07348$, $+0.17319$ and $+0.27844$. The reuse cost is $0.034$ to $0.126$
of the null, and the null is $7.95$ to $29.4$ times the reuse cost
(\texttt{real/data/rand\_control\_amd\_05\_bf16.json},
\texttt{real/data/rand\_control\_amd\_7b\_bf16.json}). So the leader's index
carries real information -- borrowing is not merely using a set the layer did
not choose -- and yet at under an eighth of the null it is far from
sufficient: most of the penalty of swapping a set is simply that the set is
not the layer's own. The 0.5B and 7B reruns reproduce this paper's published
cells exactly, with a maximum absolute difference of $0.0$ nats, so the null is
quoted against the same numbers the grid reports.

The same null on the trained-sparse checkpoint makes the point harder. There
the arbitrary set costs $+0.039778$, $+0.076422$ and $+0.102348$ nats for
$g{=}2,3,6$ against reuse costs of $-0.000686$, $+0.001180$ and $+0.002704$:
$37.8$ to $64.8$ times, against $7.95$ to $29.4$ on dense weights, and the null
exceeds reuse in every one of the nine batches
(\texttt{real/data/dsa\_group\_rand.json}). Because both cells of each pair are
measured in the same forward-pass protocol and on the same batches, the ratio is
paired rather than assembled from two experiments
(\texttt{real/data/dsa\_group\_rand\_analysis.json}). Read this as a ratio and
not as a magnitude: the trained model pays \emph{less} in absolute nats for an
arbitrary set than 0.5B does ($+0.0398$ to $+0.1023$ against $+0.3149$ to
$+0.5849$), and the two checkpoints are not the same model at two scales, so the
absolute numbers are not comparable across the two families. What is comparable
is where each model's own reuse cost sits against its own null, and on that axis
the trained checkpoint sits further below: its reuse cost is at most $2.6$\% of
its null in every cell, against $12.6$\% at worst on dense weights. So learning
a selection does not make a layer tolerant of an arbitrary one; it makes the
layer's own choice cheaper relative to the arbitrary alternative. We checked
that the two harnesses build the same null rather than two conventions wearing
one name: both exclude positions at or past the query, both
stay in range, both reproduce under a fixed seed and both change under a
different one, and both clamp identically when a row can see fewer keys than
$k$ (\texttt{real/null\_parity.py}).

\subsection{Sequence length is the amplifier}

Table~\ref{tab:long} repeats the lower half of the grid at 2048 positions.
For the same $k$ and the same group, the cost is 2.6 to 14.9 times larger.

\begin{table}[h]
\centering
\caption{Cost at 512 versus 2048 positions, same $k$ and group (0.5B,
twenty batches at both lengths). The ratio grows fastest where $k$ is far
from saturation. Source: \texttt{real/data/main\_20b\_05.json} and
\texttt{real/data/long\_20b.json}.}
\label{tab:long}
\begin{tabular}{rrrrr}
\toprule
cell & 512 & 2048 & ratio \\
\midrule
$k{=}128$, $g{=}2$ & +0.1726 & +0.7230 & 4.2 \\
$k{=}128$, $g{=}6$ & +0.3581 & +0.9254 & 2.6 \\
$k{=}256$, $g{=}2$ & +0.0242 & +0.3615 & 14.9 \\
$k{=}256$, $g{=}6$ & +0.0736 & +0.5425 & 7.4 \\
\bottomrule
\end{tabular}
\end{table}

This matters for how the reuse literature should be read. A configuration
tested at 512 positions and reported as cheap can be fifteen times more
expensive at the lengths sparse attention exists to serve, and the
amplification is strongest precisely for the large-$k$ cells that look
safest at short length: $k{=}256$ with $g{=}2$ costs $+0.0242$ nats at 512
and $+0.3615$ at 2048. The number to quote is a cost \emph{curve over
length}, not a single operating point.

\subsection{Two indices, neither of which is the axis}
\label{sec:leader}

We measure 210 single-pair transfers on 0.5B: layer $t$ borrows layer
$s$'s selection, $1 \le t - s \le 12$, all other layers left at their own
selection. Read as a marginal profile, the median cost looks nearly flat:
$+0.00067$ at distance one, $+0.00076$ at five, $+0.00116$ at ten. But that
profile pools every source that can reach a given distance, and in an
$L$-layer stack a pair at distance $d$ can only have source
$0 \ldots L-1-d$ -- so the mean source at distance $d$ is exactly
$(L-1-d)/2$, which is arithmetic rather than a property of the model. Over
the four sweeps the observed column matches that arithmetic column to two
decimals every time, so the deep columns of the profile are averages over
systematically shallower sources. Holding the source fixed instead, the cost
rises with distance on three sweeps out of four (within-source Spearman
$+0.172$ to $+0.236$, stratified permutation $p$ from $0.0057$ to $0.0609$).
Distance is therefore a real axis, not a flat one.

No single decomposition settles how large an axis it is.
Since $t = s + d$ holds for every pair, the three columns span only two
dimensions and any two of them generate the same space; whichever one is put
first absorbs the shared part. Ordering the fit by source gives the distance
increment $0.0015$ to $0.0224$ in $R^2$; ordering it by distance gives the
position increment $0.0436$ to $0.1341$; ordering it by target gives a third
set of numbers. These are three writings of one two-variable fit, not three
findings, and the near collinearity of source and target on our grids
($r = +0.82$ to $+0.99$) leaves the split unidentifiable from this data. What
survives that is the ceiling: no two-variable fit explains more than
$0.0542$ to $0.1435$ of the variance in pair cost, so $86\%$ to $95\%$ of it
is carried by something these two indices do not name. The one comparison
here that needs no ordering at all is each column fitted alone: source
explains $0.0456$ to $0.1420$ and distance $0.0012$ to $0.0268$, so the
leader is the larger axis in the weaker sense of the two marginal fits, not
in the sense that the two effects are separable from this data
(\texttt{real/data/distance\_axis.json}).

The tail is thick and not at the far end: the largest single transfer is
$+0.0635$, at distance ten (layer 11 borrowing layer 1).

Regrouping the same 210 pairs by \emph{which layer is the source} tells a
different story:

\begin{table}[h]
\centering
\caption{Mean single-pair cost by source layer, Qwen2.5-0.5B, $k=256$,
512 positions, twenty batches, distances up to twelve. Source layers 0 and 1
dominate. Source: \texttt{real/data/pairs20w12\_0.5B.jsonl}.}
\label{tab:leader}
\begin{tabular}{rrc}
\toprule
source layer & mean cost & $n$ \\
\midrule
0 & +0.012454 & 12 \\
1 & +0.010745 & 12 \\
2 & +0.003884 & 12 \\
3 & +0.000799 & 12 \\
4 & +0.000643 & 12 \\
\midrule
sources 0--1 & +0.011599 & 24 \\
sources $2+$  & +0.000972 & 186 \\
\bottomrule
\end{tabular}
\end{table}

The ratio is 11.9 per pair and 12.4 when each source layer first
contributes one mean: those are the two \emph{mean}-based weightings and they
agree. The median ratio is smaller, 5.9, because the per-pair cost has a thick
tail, so the headline describes an average rather than a typical pair. The same
split holds on the other two sizes: 18.3 and 15.9 on
1.5B, 7.4 and 6.1 on 7B. With ten thousand relabellings of the source column,
the permutation $p$ is $0.0002$ on 0.5B (the smallest model has the fewest
leader pairs) and the resolution limit $0.0001$ on 1.5B, 7B and TinyLlama
(\texttt{real/data/cross\_family.json}).

Is this a Qwen2.5 artifact? We repeat the same single-pair protocol on
TinyLlama-1.1B, a model from the Llama family, at the same $k=256$,
512 positions and twenty batches. Sources 0--1 average $+0.051198$
($n=12$) against $+0.001784$ for deeper sources, a ratio of 28.71,
with permutation $p=0.0001$ over 10000 relabellings
(\texttt{real/data/pairs20\_TinyLlama.jsonl}). On this same
window-six protocol the three Qwen2.5 sizes give 7.54, 34.99 and 11.82
(\texttt{real/data/cross\_family.json}). The leader anomaly survives a
change of model family and of scale: across the five dense checkpoints we
measure, the first two sources are the expensive ones. The fifth is
Qwen2.5-14B ($48$ layers) at $k{=}256$, where sources 0--1 average
$+0.003865$ ($n=12$) against $+0.000305$ ($n=255$) for deeper sources over the
whole $267$-pair window-six sweep (sources 2--46), a ratio of
$12.67$ with the same permutation floor $p=0.0001$; a boundary that admits
source two as "leader" dilutes it to $7.45$, so the ratio is definition-sensitive
in the direction the smaller models already show. Two things do vary: the absolute
magnitudes, and how far past the first two sources the expensive band
extends -- at $k{=}256$ and distance up to six, source two is $5.0$ times the
deeper tail on 0.5B, $22.3$ times on TinyLlama and $7.87$ times on 14B, but
indistinguishable from it on 1.5B and 7B, and the ordering \emph{within} the
first two sources is not separable at six pairs per source
(\texttt{real/data/band\_depth.json}).

That claim is scoped to dense weights, and we can say where it stops. The same
single-pair protocol on a checkpoint \emph{trained} to select its own sparse
index in every layer (a 27-layer model with a learned 2048-of-$n$ indexer,
$k{=}2048$, 4096 positions, three batches) puts lead sources at $+0.000482$
against $+0.000543$ for deeper ones, a ratio of $0.89$ with a two-sided
permutation $p=0.61$ and a difference whose bootstrap interval
$[-4.7\times10^{-4}, +4.0\times10^{-4}]$ straddles zero: on a model that
learned its own selection, the first-layer privilege is not there
(\texttt{real/data/dsa\_report.json}). Borrowing still costs something
everywhere -- $127$ of the $141$ pairs are positive, mean
$+5.4\times10^{-4}$ with interval $[+4.7\times10^{-4}, +6.1\times10^{-4}]$
-- it is simply flat in the source layer rather than concentrated in the first
two. So the asymmetry is a property of dense weights measured post hoc, not of
attention as such, and a claim that sharing is harmless because overlap is high
gains nothing from this: on that model the per-query Jaccard between shared
index sets is $0.86$ while the correlation between overlap and cost is
$0.0008$.

Grouping that model's layers costs far less than the dense grid predicts. The
dense 0.5B row at the same keep fraction ($k{=}256$ of $512$, half the sequence) pays
$+0.0242$, $+0.0114$ and $+0.0736$ nats for $g{=}2,3,6$
(Table~\ref{tab:grid}); the trained checkpoint pays $+0.000918$ on average
across the same three cells ($8$ of $9$ cells positive), which is 2.5\% of
the dense total, and each individual cell's interval includes zero at three
batches. The direction is the same, so what dense weights call the cost of
sharing is largely the cost of imposing a selection a layer never learned to
use, and the part of it that survives training is small.

These numbers are means over every scoring position, and at this keep fraction
half of those positions do not separate a borrowed set from the layer's own. A query at position $q$ sees $q+1$
keys; while $q+1 \le 256$ the top-$256$ set is every visible key, so a borrowing
layer receives exactly the set it would have chosen and that row contributes
zero. Restricting the average to the rows that can discriminate multiplies every
cell by $511/255 = 2.0039$ exactly -- measured, not fitted, and exact to
$5.4\times10^{-6}$ across twelve cells at four values of $k$
(\texttt{real/data/dilution\_analysis.json}). At $k{=}64$ the correction is only
$1.14$, so the published numbers are conservative and they are more conservative
where $k$ is a larger share of the sequence, which is the regime this paper
argues is dangerous. No comparison changes: a common factor cancels in every
ratio we report, including the null below.

\subsection{A checkpoint trained to share}
\label{sec:trainedshare}

\paragraph{The gap this leaves.} \S\ref{sec:leader} cites a checkpoint that
learned to select its own index in every layer, which is not the same as
learning to be \emph{reused}, and the Limitations section says so. A
trained-with-sharing model does exist for the architecture we study: we measure
\texttt{fxmeng/DSV2-Lite-IndexCache-Sparse-1250}, a 27-layer DeepSeek-V2-Lite
derivative whose route pattern is \texttt{FSS} repeated nine times. Each full
layer computes routes that its two shared layers then reuse; the sharing is
what the model was trained with, not something we impose on it afterwards. It
selects $2048$ of $4096$ positions.

\paragraph{What it costs to substitute the set.} We note all three readings are
paired inside one run, on the same three batches, with the same seeds
(\texttt{real/data/ic\_null.json}). Borrowing a neighbouring group's real,
learned set costs $+0.005763$ nats per token, and its interval
$[-3.2\times10^{-3}, 1.5\times10^{-2}]$ includes zero at three batches: on this
model the group's own choice is close enough to a usable substitute that this
measurement does not separate them. Replacing the same borrow with a random set
of the same size, drawn only from positions each query can see, costs
$+0.070199$ with interval $[+6.0\times10^{-2}, +8.0\times10^{-2}]$. The ratio is
$12.18$ (\texttt{real/data/ic\_analysis.json}).

\paragraph{Reading it as a ratio, not as a magnitude.} The absolute number is
smaller than on the dense checkpoints, and that is not the comparison this pair
supports: a different model family, a different keep fraction, and a different
sequence length all sit between them. What the pair does say is that training a
model to share does not make its layers indifferent to \emph{which} set they
receive. It makes the set the group learned relatively cheap; the cost of an
arbitrary substitute is still twelve times larger.

\paragraph{Two controls, and the first one lied to us.} Shuffling the order of
the borrowed indices without changing the set costs exactly $0.000$ -- as it
must, because the indices become a boolean mask and a mask is order-invariant.
We report it because our first attempt treated that zero as evidence the
intervention failed, when in fact it says only that the harness has no order
channel (\texttt{real/data/ic\_null.json}). A single-chunk probe does have to be
run where the intervention can bite, and the rule is off by one from the obvious
one: a chunk discriminates when its \emph{first} query sees more than $2048$
keys, not when its last one does. Exactly half the chunks qualify, sixteen of
thirty-two. Chunk $(1920,2048)$ sits on the boundary -- its last query sees
$2048$ keys and selects all of them -- and forcing it changes nothing
($\Delta = 0.000000$). The first genuinely live chunk $(2048,2176)$ exceeds the
cutoff by a single key and still moves by $+0.0023$; the last chunk
$(3968,4096)$, where $1921$ of $3969$ keys are excluded, moves by
$+0.0013$ to $+0.0014$
(\texttt{real/data/ic\_sanity.json}). Neither the two chunks nor the two runs
rank each other: a single-chunk substitution is a mechanism check, not a
magnitude, and this one moves by $2.7\%$ between separate processes while
restoring bit-identically to $0.000000$ inside its own run. What it establishes
is that the intervention reaches the model at all, which the first attempt --
aimed at a chunk that could not discriminate -- did not.

\paragraph{Half the positions are inert here too.} The dilution of
\S\ref{sec:leader} reproduces on this checkpoint, and here it is measured rather
than derived. Applying the substitution to every chunk and to only the sixteen
discriminating chunks gives the same three per-batch deltas to the last printed
digit (\texttt{real/data/ic\_dilution.json}). Averaging over the $2047$ rows that
can discriminate instead of all $4095$ multiplies the cells by
$1.9932$, $1.9913$ and $2.0136$ against the predicted $4095/2047 = 2.0005$
(\texttt{real/data/ic\_row\_dilution.json}).

\paragraph{Overlap still does not carry the signal.} On the last chunk, where
both sets can exclude half the sequence, the learned adjacent-group set overlaps
a random set of the same size at Jaccard $0.5211$ against $0.3399$ -- a factor
$1.53$ -- and yet the two differ in cost by the factor above
(\texttt{real/data/ic\_overlap.json}). Random sets have to be compared inside
the same chunk, because their own overlap falls with the visible range: from
$0.9418$ at the first chunk to $0.3399$ at the last.

\subsection{What layer-specific scoring buys}
\label{sec:latentindex}

The checkpoint above gives every follower the leader's score. A second
checkpoint from the same family, \texttt{fxmeng/DSV2-Lite-LatentIndex-Sparse-1250},
uses the same $FSS$ route pattern but a different rule: the leader shares only
its latent cache, and each follower reads its own offset out of the route list
and scores the same keys with its own projection (the construction
\citet{latentindex} describes). That makes the two checkpoints a controlled
comparison of one thing -- whether the follower scores with its own weights or
with the leader's -- so we can ask what that rule buys.

\paragraph{Nothing measurable, and the probe did fire.} Forcing every follower
to read offset zero reproduces the leader-scored behaviour, and a repeat with
the substitution applied on top of each scores the cost of an arbitrary set:

\begin{center}
\begin{tabular}{lc}
\toprule
cell & mean $\Delta$nll \\
\midrule
own score, own set (control) & $+0.000000$ \\
leader score, own set & $+0.000105$ \\
own score, random set & $+0.069073$ \\
leader score, random set & $+0.069093$ \\
\bottomrule
\end{tabular}
\end{center}

The difference between the two ``own set'' cells has a confidence interval
$[-0.001347,\,0.001556]$ over the three batches, which contains zero, and the
three per-batch values do not even share a sign. The ratio between the two
random-set cells is $0.9996$ by mean of per-batch ratios (a ratio of means
gives $0.9997$), and the substitution cost is $0.984$ times the cost measured
on the leader-scored checkpoint
(\texttt{real/data/li\_run.json}, \texttt{real/data/li\_analysis.json}).

A ratio near one has two possible readings -- the quantity is flat, or the
probe did not reach the model -- and only a separate check separates them.
Re-capturing the selected indices with the offset forced shows that they do
change: on the $306$ chunks where the first query can see more than $k$ keys,
Jaccard between the own-score and leader-score selections averages $0.7799$
(median $0.7878$, minimum $0.4322$), and not one chunk is left identical. On
the $288$ chunks that do not discriminate, $97.6\%$ are identical, which is
what should happen there: their scores are all $-\infty$ and top-$k$ is a
tie-break rather than a decision. The measurement is flat because the quantity
is flat.

\paragraph{What this rules out.} It rules out the inference that better
selection quality must mean a lower decision-level cost. Giving a follower its
own score changes roughly a fifth of the set it selects and leaves the cost of
an arbitrary substitute unchanged. It does not rule out that layer-specific
scoring helps on the axis its authors measure -- attention-mass recall
against the dense attention is not the same quantity as NLL after
substitution, and a method can recover one without moving the other. We
therefore state the first and not the second, having been rejected once for
making exactly that inference.

\subsection{A predictive diagnostic, and a construct check}
\label{sec:mechanism}

The split by source says \emph{who} is expensive, not why. The obvious
explanations fail: overlap, linear CKA and the contribution proxy do not rank
the cost (\S\ref{sec:whatnot}). We therefore explore one further quantity, and
we stress at the outset that it is a \emph{predictive diagnostic}, not an
identified mechanism.

The quantity is where a leader's selection lands in the \emph{follower's}
distribution. We read the leader's selected positions against the follower's
dense softmax (no reuse mask) and sum their probability; call this
\emph{deliver-mass}. This is an unnormalised variant of the attention-mass
overlap ratio used as a systems heuristic by Kascade \citep{kascade} and of the
attention-mass recall reported by LatentIndex \citep{latentindex}; we apply it
to dense-weight post-hoc reuse rather than claim it as a new measure. On 0.5B,
sources 0--1 deliver $0.9087$ against $0.9722$ for deeper sources, a ratio of
$0.935$; on TinyLlama the split is $0.8352$ versus $0.9784$, a ratio of $0.854$
(both $p=0.0001$, ten-thousand relabellings). The leaders' own selections stay
high-mass under their own distributions, so the gap is a cross-layer mismatch
rather than diffuse early attention; Kascade \citep{kascade} independently
observes that early layers have flatter distributions.

Across the paired probes deliver-mass correlates with reuse cost at $-0.5124$
on 0.5B and $-0.6912$ on TinyLlama. Deliver-mass also rises with layer depth,
so we partial out the continuous source index: the correlation stays $-0.4679$
and $-0.6290$. This control is incomplete: the leader reset is really the
\emph{binary} split sources $0$--$1$ versus $2+$, and controlling that dummy
shrinks the correlation materially: to $-0.3338$ ($p=1.6\times10^{-4}$) on
0.5B and $-0.4253$ ($p=3.3\times10^{-6}$) on TinyLlama, while deliver-mass
still adds $+0.0915$ and $+0.1149$ explained variance beyond the dummy alone,
and a stack controlling the dummy, the source index and the distance together
leaves $-0.3650$ and $-0.4197$ (\texttt{real/data/dummy\_controls.json},
battery in \texttt{real/dummy\_controls.py}). Deliver-mass is therefore partly
entangled with leader identity but is not a restatement of it. Restricting to
sources $2+$ keeps
TinyLlama negative at $-0.7108$ but weakens 0.5B to $-0.2957$; shuffling the
deliver labels destroys the signal ($-0.0911$ and $-0.0431$). The correlations
and the shuffle control are read from \texttt{real/data/deliver\_corr.json},
joined from the dense probe and the masked-NLL pairs by
\texttt{real/leader\_deliver\_analysis.py}. A reviewer asked whether this is
just a decorative restatement of plain selection overlap. We tested that
directly: on the \emph{same} dense pass we computed the Jaccard overlap of the
two top-$k$ index sets and compared the two predictors on the same pairs
(\texttt{real/overlap\_head2head.py}, \texttt{real/data/h2h\_corr.json}). Jaccard
does carry signal ($r=-0.363$ on 0.5B, $-0.345$ on TinyLlama, $-0.427$ on 1.5B
and $-0.402$ on 7B), but deliver-mass
keeps a strong partial correlation after controlling it ($-0.4218$ and
$-0.6469$ at the two small models, $-0.631$ and $-0.625$ at 1.5B and 7B,
$p\le1.2\times10^{-6}$) and adds variance ($\Delta R^2=+0.154$, $+0.369$,
$+0.326$ and $+0.328$), while jaccard's partial collapses and even flips
positive on
TinyLlama ($+0.138$, $p=0.149$). A synthetic fixture where deliver is a pure
monotone function of overlap makes both the partial and $\Delta R^2$ vanish, so
the comparison can go red. Under this test deliver-mass is not redundant with
plain overlap. This is still a predictive comparison, not a mechanism: both
readouts share the follower's queries and keys with the masked-NLL cost, and
we note that we make no causal claim.

The binary leader control of the previous paragraph is scale-dependent, and we
report that plainly rather than only where it is strongest
(\texttt{real/data/dummy\_controls.json}). At 1.5B the leader flag already
explains half the cost variance ($R^2=0.498$), so deliver-mass adds only
$+0.028$ beyond it and the partial falls to $-0.237$ ($p=3.9\times10^{-3}$);
at 7B the flag explains $0.306$ and deliver adds $+0.166$ with a partial of
$-0.489$ ($p=3.4\times10^{-10}$). The 0.5B and TinyLlama readings sit between
these ($+0.092$ and $+0.115$ added variance). The correlation survives the
binary control at every scale, but at 1.5B it does so weakly, and a reader
should not conclude that the diagnostic is uniformly useful across sizes.

A correlation surviving controls is not the same as accounting for the size of
the effect, and we note the two are far apart here. Fitting cost against
deliver-mass on the late leaders alone and extrapolating that slope to the early
leaders' mean deliver predicts only $32\%$, $60\%$, $28\%$ and $38\%$ of the
observed early-minus-late cost gap at 0.5B, TinyLlama, 1.5B and 7B
(\texttt{real/data/leader\_magnitude.json}, battery in
\texttt{real/leader\_magnitude.py}); the remaining $40$--$72\%$ is unexplained by
the mass shift even under this deliberately generous linear extrapolation. The
extrapolation is also out of support: fewer than half the early pairs fall inside
the late group's deliver range, and at 1.5B none do (early maximum $0.9242$ below
late minimum $0.9379$), so deliver-mass is nearly collinear with the leader flag
there. We therefore report deliver-mass as a correlated diagnostic that falls
well short of explaining why the first layers are expensive, not as an account of
that magnitude; what carries the residual is not isolated by this measurement.

As a construct check we ask only whether moving the borrowed index moves the
cost in the predicted direction. A pre-registered two-pass ORACLE
(source: \texttt{real/reuse\_intervene.py}) reads dense mass once and replaces
only the index followers borrow, holding each layer's $q/k/v$ and the leaders'
own selections fixed. HIGH borrows the followers' consensus dense top-$k$,
LOW the bottom-$k$; BASE is the leader's own selection. On the same bfloat16
0.5B artifact as the grid, HIGH$<$BASE$<$LOW in every grouping (group 2/3/6
BASE NLL $2.874$ / $2.874$ / $2.911$, LOW $4.26$--$5.23$;
\texttt{real/data/intervene\_05\_v2.json}), and the ordering holds on
TinyLlama. We note this ordering is partly an artifact of construction: HIGH
and LOW are the top and bottom of the same follower dense distribution that
defines deliver-mass. The check therefore shows the borrowed index is a live lever on
cost, \emph{not} that deliver-mass identifies a mechanism, and the guide set is
dense attention unavailable in an autoregressive decoder, so this is not a
deployable selection rule. We further note the ORACLE's BASE index is not
bit-identical to the grid's reuse condition (a same-batch diagnostic,
\texttt{real/diag\_base.py}, finds a material NLL gap and selection positions
that differ at most indices), so we make no cross-comparison of absolute
magnitudes between the two experiments.

One descriptive consequence on 0.5B is the width non-monotonicity. In an even
grouping layer 3 is the follower of layer 2, and that single transfer costs
$+0.014557$, 18.9 times the median of all 210 pairs
(\texttt{real/data/pairs20w12\_0.5B.jsonl}). Removing layer 3 from the
$g{=}2$ partition drops its sum from $+0.023558$ to $+0.009001$, a factor of
2.6. Layer 3 is also the dearest follower of layer 0 ($+0.036897$) and the
third dearest of layer 1 ($+0.006638$, against $+0.000231$ for layer 2).
It is not an odd-versus-even effect: odd followers average $+0.001963$ and
even ones $+0.000569$, but removing layer 3 turns the odd average into
$+0.000818$, close to the even one. We stress this is a single-layer
observation on one width: on 1.5B and 7B the analogous layer is not the
outlier, so we do not offer a width-dependent mechanism. Groups of three do
not ask layer 3 to follow, which is why the offset fix helps.

Second, the fix is cheap but not uniform. Shifting every group by one
layer so the leaders become layers 1, 4, 7, \ldots{} multiplies the cost by
0.30, 1.53, 0.31 and 0.50 for $g{=}2,3,4,6$
(\texttt{real/data/offset20.jsonl}). Three widths get 2--3.4 times cheaper;
$g{=}3$ gets 1.5 times dearer, and the reason is the same layer 3. In the
unshifted $g{=}3$ partition layer 3 leads and pays nothing; shifting by one
makes it a follower of layer 1, and the partition's predicted sum rises from
$+0.010182$ to $+0.017952$, which matches the measured $+0.017464$. A system
that shares routing can choose which layer computes it, and the choice is
worth a factor of three \emph{in the partitions that do not put layer 3 in a
follower slot}.

\subsection{Borrowed costs add up}

If the single-pair costs were independent, a group's cost would be the sum
of its members'. It is not, and the departure grows with group width. On
0.5B, measured over summed-predicted is 1.03 at $g{=}2$, 1.12 at $g{=}3$,
1.16 at $g{=}4$, 1.32 at $g{=}6$ and 2.90 at $g{=}12$: $g{=}2$ is additive
to within three percent, while the widest grouping we ran is nearly three
times the sum of its parts. On 1.5B the sum understates the group cost
throughout, by a factor of 1.72 to 2.32; on 7B it overstates it, by a factor
of 0.40 to 0.71. The sign of the departure therefore changes with model
size, so we report a per-model factor rather than additivity as a property
of the model (\texttt{real/data/additivity.json}).

What does transfer between the two protocols is the \emph{ordering}. Across
the five widths on 0.5B and the three on each other model, the rank
correlation between the summed single-pair cost and the measured group cost
is $+1.000$ in all three. So the single-pair protocol is a usable instrument
for asking which grouping is cheaper and an unusable one for asking how much
cheaper. The direction of the error is not even constant: on 7B the summed
single-pair cost overstates the group cost at every width, by up to a factor
of two and a half, while on 0.5B it understates the widest width by a factor
of three.

The widest cell is also the one where an explanation is available but not
separable from the measurement. At $g{=}12$ the 24 layers form two groups,
so eleven consecutive layers borrow from layer 0 at once and each one's
corrupted output is the next one's input; at $g{=}2$ no follower of layer 0
has another follower of layer 0 upstream of it. Chain depth and
group width move together in this grid and we did not run the design that
would separate them, so we record the pattern and not a decomposition of
it.

\section{What the similarity measures do and do not predict}
\label{sec:whatnot}

The reuse literature's premise is that overlap is high, therefore reuse is
cheap. We measured the cost and the overlap on the same models, so the
implication can be tested rather than assumed. The three checks below come from
two different corpora, which we keep separate: overlap is on 23 Qwen2.5-0.5B
layer pairs, while the CKA and contribution probes are an earlier two-text
GPT-2 study (eleven and ten layer pairs). They agree in sign but are not a
single controlled comparison.

\paragraph{Overlap.} We recount overlap per (batch, head, query), comparing
each layer's selection with its neighbour's, and normalise by the overlap
expected if the two selections were independent given each query's prefix
length. That expectation is not $k$ divided by the prefix: for a prefix shorter
than $k$ the selection is the whole prefix, so a naive baseline saturates at
one when $k$ approaches the sequence length and manufactures a negative effect
at exactly the cells where reuse is cheapest. The correct closed form is
$T^{-1}\sum_{u \le T} \min(k, u)/u$. With it, enrichment is 1.19 times chance
at 512 positions with $k=128$, 1.03 at $k=256$, 1.96 and 1.46 at 2048
(\texttt{real/data/overlap\_baseline.json}).

Across 23 layer pairs at matched $k$, overlap and cost have $r = -0.357$
(rank correlation $-0.453$), and the overlap range is narrow, 0.83 to 0.93.
The sign is the one the premise predicts -- more overlap, cheaper reuse -- so
what fails is the magnitude and the exclusivity, not the direction: over this
narrow overlap band the cost spans an order of magnitude, and the relation is
not monotone in the way a predictor would need to be.

\paragraph{Representation similarity.} Linear-kernel CKA between adjacent layers'
attention outputs gives $r = -0.341$ with cost over ten layer pairs after
excluding the first. Note that linear CKA is biased upward in finite samples:
over twenty seeded draws of two independent $n\times 64$ matrices it averages
$0.243$ at 200 samples and $0.031$ at 2000, so a value is only interpretable
next to $(n,d)$ and we report both with every number.

\paragraph{Contribution magnitude.} A layer's attention output, pushed through
the unembedding, has a norm that grows 33-fold from layer one to layer eleven.
It is strongly correlated with depth ($+0.895$ on the log scale) and with cost
at $r=-0.071$ overall, $+0.604$ after excluding the first pair. Our
preregistered criteria required an unrelated-quantity control to correlate
less than 0.5 in absolute value with depth; it fails at 0.895, so this proxy
is rejected rather than reported as weak.

None of the three is a sufficient predictor, and each falls short for an
interpretable reason: overlap is a property of attention rather than of a
layer's role, and much of its apparent signal is shared with leader identity;
CKA measures similarity of outputs, not of what those outputs are needed for;
contribution magnitude is a depth proxy. We report them because the natural
predictor is weaker and less exclusive than the premise assumes, which is what
leaves room for the mass diagnostic we turn to next.

\section{Decision-level cost}

\subsection{Power first}

A token-level cost of 0.06 nats and an unchanged answer can coexist, and
reporting only the token-level number would overstate the practical cost.
We therefore measure accuracy, but we measure the \emph{detectable} accuracy
difference at the same time, because a small sample with few discordant
items is unable to distinguish ``no damage'' from ``damage below the resolution''.

For each cell with $b$ answers changed from right to wrong and $c$ from wrong
to right, the exact two-sided McNemar test is applied, and the minimum
detectable difference is the smallest $|b-c|$ that would reach $p<0.05$ given
the observed discordance count. Table~\ref{tab:mde} reports the ranges.

\begin{table}[h]
\centering
\caption{Minimum detectable accuracy difference, in percentage points, for
cells that admit a rejection region at all. Source:
\texttt{real/data/mcq\_power.json}.}
\label{tab:mde}
\begin{tabular}{llrr}
\toprule
benchmark & model & cells with MDE & magnitude \\
\midrule
ARC-Easy & 0.5B & 20 & 3.0--9.5 pt \\
ARC-Easy & 1.5B & 5 & 3.5--7.0 pt \\
ARC-Easy & 7B & 2 & 4.7 pt \\
\bottomrule
\end{tabular}
\end{table}

Cells whose discordance count is too small to reject anything are labelled
as having no rejection region rather than given an MDE: at two discordant
items the smallest two-sided exact $p$ is 0.5, so no result could reach
significance. Giving those cells a number would be an arithmetic fiction.

\subsection{The token-level cost does not translate one-for-one}

On ARC-Easy with 0.5B, 25 cells, dense accuracy 0.6100:

\begin{itemize}
\item Only five cells reach $p<0.05$, and all of them have a group of six
      layers or more: $k{=}32$ with $g{=}12$ ($-25.5$ points), $k{=}64$ with
      $g{=}12$ ($-21.5$), $k{=}128$ with $g{=}12$ ($-21.0$), $k{=}32$ with
      $g{=}6$ ($-9.0$), $k{=}192$ with $g{=}12$ ($-8.5$).
\item $k{=}256$ is flat: five cells, five discordant items in total, every
      $p \ge 0.5$, and two cells with exactly zero.
\item The maximum difference across all 25 cells is 25.5 points, and all
      five significant cells have it; among the twenty cells that do not
      reach significance the largest is 4.5. At $k{=}256$ it is 1.0.
\end{itemize}

The 1.5B grid repeats the shape more quietly: dense 0.7600, nine cells, all
three $k{=}256$ cells with zero discordant items, and a largest difference of
13.0 points (significant). The 7B grid is quieter still: dense 0.7933, six
cells, largest difference 4.7 points, and every $k{=}256$ cell at 0.67 from a
single wrong-to-right flip.

\subsection{Banding}

Sorting the 25 cells of 0.5B by token-level cost reveals where the translation
starts:

\begin{table}[h]
\centering
\caption{Accuracy change by token-level cost band, 0.5B, 25 cells. Source:
\texttt{real/data/mcq\_summary.json} crossed with
\texttt{real/data/main\_20b\_05.json}.}
\label{tab:band}
\begin{tabular}{rrrrr}
\toprule
cost (nats) & cells & mean change & worst change & smallest $p$ \\
\midrule
0.0--0.2  & 11 & $-0.7$ pt & 2.0 & 0.219 \\
0.2--0.5  &  5 & $-1.8$ & 3.0 & 0.26 \\
0.5--1.0  &  5 & $-4.7$ & 8.5 & 0.0076 \\
1.0--2.0  &  2 & $-15.0$ & 21.0 & $4.3\times10^{-7}$ \\
2.0--4.0  &  2 & $-23.5$ & 25.5 & $1.7\times10^{-9}$ \\
\bottomrule
\end{tabular}
\end{table}

Pearson correlation between the two is $-0.936$ over 25 cells, and $-0.906$
over nine cells on 1.5B, $-0.638$ over six on 7B. The correlation is not
evidence that NLL substitutes for accuracy (cells share question sets and
are not independent), but it rules out the worse reading: it is not true
that token-level damage is uniformly harmless downstream.

The defensible statement is banded. Below 0.2 nats, we detect no decision
damage at these sample sizes. Above 0.5 nats, we do.

\subsection{Why the answers mostly survive}

A common offset that moves all four option scores by the same amount changes
no argument. Decomposing each item's score change into a mean shift across
options and a spread across options separates the two:

\begin{table}[h]
\centering
\caption{Score decomposition at 0.5B, 120 items, $g{=}6$. The shift is
large relative to the baseline margin, so the reason answers survive is
that margins were wide, not that the perturbation was small. Source:
\texttt{real/data/mcq\_offset\_05b\_k64g6.json} and its $k{=}128$ counterpart.}
\label{tab:offset}
\begin{tabular}{lrr}
\toprule
quantity & $k{=}64$ & $k{=}128$ \\
\midrule
items flipped & 27 / 120 & 25 / 120 \\
mean shift & $-2.541$ & $-1.077$ \\
mean spread & 4.376 & 2.953 \\
baseline margin & 2.832 & 2.887 \\
AUC(margin predicts flip) & 0.764 & 0.856 \\
AUC(spread predicts flip) & 0.616 & 0.579 \\
AUC(shift predicts flip) & 0.514 & 0.648 \\
margin, flipped & 1.201 & 0.746 \\
margin, not flipped & 3.306 & 3.450 \\
\bottomrule
\end{tabular}
\end{table}

The mechanism is not that the perturbation is small: the spread of 4.376 exceeds
the baseline margin of 2.832, so the ordering of options is genuinely at
risk. It survives because the margin \emph{of the items that flipped} was
already thin ($1.20$ against $3.31$), and margin predicts which items flip
with an AUC of 0.76 to 0.86. A system whose answers are near-ties will lose
them; one whose answers are decided by a wide margin will not. Neither the
common shift nor the spread predicts as well as the pre-existing margin does.

\subsection{Direction, and a claim we had to retire}

Our first report said the flips were directional at 5.1 to 1 on 1.5B. That
claim was wrong, and the reason is algebraic:

\begin{equation}
b - c = (\text{correct under own}) - (\text{correct under reuse})
      = N\,(\text{own accuracy} - \text{reuse accuracy}).
\end{equation}

This holds exactly on all 40 cells we checked (residual 0.000), and we now
assert it in the analysis script rather than in prose. Dividing by $b+c$
therefore restates the accuracy drop in another denominator; on a cell whose
accuracy is preserved, $b$ and $c$ are equal as an algebraic consequence. The
ratio is not evidence of direction.

What can carry direction is a pair of \emph{different} denominators, because
$b$ can only come from items originally correct and $c$ only from items
originally wrong:

\begin{equation}
\text{damage} = \frac{b}{\#\text{correct under own}}, \qquad
\text{rescue} = \frac{c}{\#\text{wrong under own}}.
\end{equation}

\begin{table}[h]
\centering
\caption{Damage and rescue per hundred at-risk items. ARC-Challenge's at-risk
population is \emph{inverted} (more wrong than correct), which rules out
drift-toward-the-majority as the explanation. Source:
\texttt{real/data/cross\_bench.json}.}
\label{tab:atrisk}
\begin{tabular}{lrrrr}
\toprule
benchmark & at risk (correct / wrong) & damage & rescue & asymmetry \\
\midrule
ARC-Easy & 5098 / 2602 & 10.1 [9.3, 11.0] & 7.9 [6.9, 9.0] & 1.28 \\
ARC-Challenge & 1245 / 1275 & 12.4 [10.7, 14.4] & 6.9 [5.6, 8.4] & 1.80 \\
\bottomrule
\end{tabular}
\end{table}

Both benchmarks agree in direction and disagree in magnitude by only a
factor of 1.4. The ARC-Challenge result is the stronger one for our
interpretation, because its at-risk population is inverted: drift toward
whichever class is larger would flip the sign, and the sign does not flip.

\subsection{A confound we did not remove: exemplar count}

The prompt builder discards demonstrations rather than truncating the
question, so the number that fit depends on question length. ARC-Easy fits
all eight; ARC-Challenge fits five. Comparing the pooled rates across
benchmarks therefore changes two things at once.

\begin{table}[h]
\centering
\caption{The same six cells at two exemplar counts (0.5B, ARC-Easy,
$k \in \{64,128\}$, $g \in \{2,3,6\}$). The sign of damage minus rescue
reverses. Source: \texttt{real/data/cross\_bench.json} under
\texttt{shot\_control}.}
\label{tab:shots}
\begin{tabular}{rrrrrr}
\toprule
exemplars & $b$ & $c$ & damage & rescue & asymmetry \\
\midrule
8 & 74 & 54 & 10.3 [8.2, 12.7] & 11.3 [8.6, 14.4] & 0.91 \\
5 & 45 & 24 & 6.1 [4.5, 8.1] & 5.2 [3.4, 7.6] & 1.17 \\
\bottomrule
\end{tabular}
\end{table}

The reversal is asserted in the instrument's self-test, so the caveat is
machine-checked rather than promised. Two consequences: the gap between
1.28 and 1.80 is not attributable to question difficulty alone, since the
matched subset moves the other way; and exemplar count is a real variable
in this measurement, not noise. To align it would require either eight exemplars
at a larger budget, which changes the absolute number of reusable keys, or
five-shot ARC-Challenge, which we did not run.

\section{Limitations}
\label{sec:limits}

\paragraph{Post hoc sharing.} Most numbers here borrow routing from a model
trained without sharing. \S\ref{sec:leader} adds a checkpoint that learned to
select its own sparse index, and \S\ref{sec:trainedshare} adds one trained with
sharing (an FSS route pattern, each full layer serving two shared ones). That
closes the gap in the direction a reader would ask for, and leaves two smaller
ones open. It measures substitution within one architecture derivative at one
keep fraction, not the sweep a deployed system would need; and it can only say
that the group's learned set is a cheap substitute, since the interval on that
reading includes zero at three batches. A model trained to share could also move
its representations so that sharing becomes cheaper still, which no post-hoc
substitution can observe -- we compare sets handed to frozen weights, not a
model that learned to receive them.

\paragraph{No downstream generation.} We score multiple choice by summed
option log-probability and do not generate. Free-form generation would change
the attribution, but it is also a setting we have not measured.

\paragraph{Four-option benchmark resolution.} At 200 items, an MDE of 3--9
points is coarse. Cells reported as indistinguishable are that, not
demonstrated equivalence.

\paragraph{Three model sizes, three data regimes.} The grid covers three
sizes because that is what we ran it on: 0.5B has all thirty cells, 1.5B and
7B have nine each, and 14B has three. The 7B token-level cells exist but only
for the three $k$ values that fit, with twenty batches and bfloat16; its
negative deltas at four batches turned positive at twenty, so we do not have a
size at which the sign of the group cost is stable to sample size. The leader
protocol reaches 14B completely: the window-six sweep over $48$ layers is
$267$ pairs, sources 0--46, and all $267$ were measured. So the anomaly is
confirmed at the top of the size range over the whole sweep, not a truncated
band of it.

\paragraph{Harness is the instrument.} Three defects in our own implementation
produced plausible numbers before the control caught them. The control is
reported in every artifact for that reason; a reader who distrusts the
numbers should distrust the harness first, which is what we did.

\paragraph{The mass diagnostic is not an identified mechanism.} Deliver-mass
is measured on dense attention while the cost is measured after masking, so
the two share the follower's queries and keys; controlling the binary
source-$\le 1$ split shrinks the correlation by about a third. We note the ORACLE
intervention brackets cost by construction (its HIGH and LOW are the top and
bottom of the same follower dense distribution), its BASE index is not
bit-identical to the grid's reuse condition, and its guide set requires dense
attention unavailable during autoregressive decoding. Two further limits are
visible in the numbers themselves: the incremental value of the diagnostic is
scale-dependent (on 1.5B the binary leader flag already accounts for half the
variance and deliver-mass adds little), and a slope extrapolation from the late
layers leaves 40--72\% of the leader cost gap unexplained. We note these
readings are a predictive diagnostic and a construct check, not a causal
account or a deployable remedy.

\paragraph{Statistical units and post-hoc power.} The leader permutation
treats single pairs as exchangeable although source and target depths are
coupled; the damage-versus-repair confidence intervals treat paired counts as
independent; the four-batch and twenty-batch results are nested; and the MDE
is computed after seeing the data. These are caveats on the decision-level
analysis, not on the token-level NLL measurements.

\section{What is not established here}

We list the statements we considered and rejected, because each was tempting
and each is unsupported:

\begin{itemize}
\item That reuse is free at large $k$. It is free at 512 positions with
      $k$ at 75\% of the length and small groups; at 2048 positions
      $k$ equal to 25\% of the length already costs $+0.1093$ nats.
\item That overlap between layers is the reason reuse works. Overlap is high,
      and its correlation with cost is negative and weak.
\item That deliver-mass identifies a \emph{mechanism}. It improves on plain
      overlap as a predictor in a head-to-head test (\S\ref{sec:mechanism}),
      but both readouts share the follower's queries and keys with the cost,
      so we note the predictive gain is not a causal identification.
\item That layer 3 is a width-dependent mechanism. The 18.9-fold outlier is a
      single-layer observation on 0.5B and does not recur at 1.5B or 7B.
\item That the direction of damage is strongly asymmetric. It is asymmetric
      at 1.3--1.8 times, and one control flips its sign.
\item That MMLU would replicate the decision-level result. All three models
      score at or near chance on it in our 5-shot setup, so a perturbation
      there measures the floor rather than the cost. We dropped it and recorded
      why.
\item That ARC-Challenge is harder for reuse because its questions are harder.
      Its exemplar count differs from ARC-Easy's, and we did not separate the
      two.
\end{itemize}

\section{Related work}

\paragraph{Cross-layer index reuse: throughput-oriented.} Multi-layer Indexer
\citep{mli} reports that neighbouring layers select nearly identical tokens
and reuses the index without retraining, targeting decode latency. CLSA
\citep{clsa} shares routing and key-value blocks to cut prefill and storage.
TidalDecode \citep{tidaldecode} makes sparse indices position-persistent across
decoding steps and HySparse \citep{hysparse} combines oracle token selection
with KV sharing. These lines measure throughput and report quality as retained.

\paragraph{Cross-layer index reuse: quality measured.} IndexCache
\citep{indexcache} is the closest prior measurement: it reuses indices across
layers in a trained sparse model and reports the quality cost at benchmark
granularity by retention ratio, including a training-free variant that drops
without search and a study of why layer similarity is a poor proxy. Kascade
\citep{kascade} is a training-free system that selects exact top-$k$ indices
in anchor layers and reuses them between, choosing anchors by a
cross-layer-similarity objective; its similarity ratio is the normalised form
of the mass diagnostic we use in \S\ref{sec:mechanism}, and it too observes
that early layers have flatter distributions. LatentIndex \citep{latentindex}
starts from a shared latent cache but gives each follower a layer-specific
scorer precisely because a shared discrete selection degrades quality, and
reports attention-mass recall. Our contribution is not the first measurement
of the quality term: it is a controlled isolation of the selection-set
variable on dense pretrained weights, holding $q/k/v$ and batch fixed with a
three-readout design, and a separation of token-level cost from its
detectability at the decision level.

\paragraph{Key-value sharing.} A systematic study \citep{kvshare} evaluates
cross-layer key-value configurations that train the shared model, reporting
that many stay competitive at half the cache. KVSharer \citep{kvsharer} finds
that sharing \emph{dissimilar} KV caches better preserves performance — an
early statement that layer similarity is a poor guide to sharing, which our
three failed predictors reproduce and extend at the token level. Trained
sharing can adapt representations; post-hoc routing has no such adaptation.

\paragraph{Measuring what a component is needed for.} Our predictive proxies
are all similarity-shaped, and all fail, consistent with the KV-sharing and
index-reuse results above. More broadly, representational similarity is not
necessity: linear CKA \citep{cka} measures alignment, not whether a component's
computation is used, and work on attention sinks \citep{streamingllm} shows
that a few positions can be load-bearing for reasons that top-$k$ saliency
alone does not capture. These are the limitations we inherit when we read a
mass-based diagnostic as a predictor.

\section{Reproducibility}

Each number in this paper is recomputed from a stored artifact by two
independent gates that run in the repository's test path, and a stale literal
would fail both: one checks that
each literal in the prose matches a value recomputed from the artifacts, in
both directions (a stale literal fails it, and a literal without an artifact
behind it fails it too);
another checks that each quoted sentence appears verbatim in its report and
that its numbers match a named artifact field within a tolerance. The
evaluation-evidence export lists 331 episodes with per-field provenance,
marks three wall-clock fields as unavailable rather than estimating them,
and refuses to cite any episode whose control did not pass. A large language
model assisted with drafting and editing; all measurements, scripts,
and reported numbers were produced and verified by the checks named above.

\bibliography{main}
\bibliographystyle{tmlr}

\end{document}
